\section{Related Work}
We organize related work into four categories: agent security benchmarks, multi-agent security research, guardrail systems for LLMs and agents, and software engineering benchmarks. We position \proposed~ relative to each category and highlight the novel contribution of measuring harm amplification in multi-agent orchestration.

\begin{table*}[!h]
\caption{Positioning of \proposed~ relative to related agent-security works. The comparison focuses on capabilities required to study local-to-global harm amplification in multi-agent LLM systems. 
\cmark{} indicates explicit support, \pmark{} indicates partial support, and \xmark{} indicates lack of primary focus.}
\label{tab:positioning}
\centering
\scriptsize
\setlength{\tabcolsep}{0.7pt}
\renewcommand{\arraystretch}{1}
\begin{tabular}{lccccccccc}
\toprule
Method / Benchmark
& \makecell{Multi-agent\\orchestration}
& \makecell{Compromised\\agents}
& \makecell{Collusion\\axis}
& \makecell{Memory\\persistence}
& \makecell{Trajectory /\\trace scoring}
& \makecell{Harm / severity\\scoring}
& \makecell{Defense\\evaluation}
& \makecell{Utility / cost\\trade-off}
& \makecell{Amplification\\metric} \\
\midrule

\multicolumn{10}{l}{\textit{Agent-security and agent-evaluation benchmarks}} \\
ASB~\cite{zhang2025agent}
& \xmark & \pmark & \xmark & \cmark & \pmark & \pmark & \cmark & \pmark & \xmark \\

AgentDojo~\cite{debenedetti2024agentdojo}
& \xmark & \xmark & \xmark & \xmark & \pmark & \pmark & \cmark & \cmark & \xmark \\

AgentHarm~\cite{andriushchenko2025agentharm}
& \xmark & \xmark & \xmark & \xmark & \xmark & \cmark & \pmark & \pmark & \xmark \\

SEC-bench~\cite{lee2026secbench}
& \pmark & \xmark & \xmark & \xmark & \pmark & \pmark & \pmark & \pmark & \xmark \\

SafePro~\cite{zhou2026safepro}
& \xmark & \xmark & \xmark & \xmark & \pmark & \cmark & \pmark & \pmark & \xmark \\

SWE-bench~\cite{jimenez2024swebench}
& \xmark & \xmark & \xmark & \xmark & \xmark & \xmark & \xmark & \xmark & \xmark \\

\midrule
\multicolumn{10}{l}{\textit{Multi-agent adversarial and persistence benchmarks}} \\
TAMAS~\cite{kavathekar2025tamas}
& \cmark & \cmark & \cmark & \xmark & \cmark & \pmark & \xmark & \pmark & \xmark \\

Who is the Mole?~\cite{xie2025s}
& \cmark & \cmark & \pmark & \xmark & \pmark & \pmark & \cmark & \pmark & \xmark \\

Memory poisoning~\cite{sunil2026memory}
& \pmark & \pmark & \xmark & \cmark & \pmark & \pmark & \cmark & \pmark & \xmark \\

\midrule
\multicolumn{10}{l}{\textit{Guardrail and monitoring methods}} \\
LlamaGuard~\cite{inan2023llama}
& \xmark & \xmark & \xmark & \xmark & \xmark & \cmark & \cmark & \pmark & \xmark \\

Qwen3Guard~\cite{zhao2025qwen3guard}
& \xmark & \xmark & \xmark & \xmark & \xmark & \cmark & \cmark & \pmark & \xmark \\

ShieldGemma~\cite{zeng2024shieldgemma}
& \xmark & \xmark & \xmark & \xmark & \xmark & \cmark & \cmark & \pmark & \xmark \\

WildGuard~\cite{han2024wildguard}
& \xmark & \xmark & \xmark & \xmark & \xmark & \cmark & \cmark & \pmark & \xmark \\

LlamaFirewall~\cite{chennabasappa2025llamafirewall}
& \pmark & \pmark & \xmark & \xmark & \pmark & \pmark & \cmark & \pmark & \xmark \\

ToolSafe~\cite{mou2026toolsafe}
& \pmark & \pmark & \xmark & \xmark & \pmark & \pmark & \cmark & \pmark & \xmark \\


AgentAuditor~\cite{luo2026agentauditor}
& \pmark & \pmark & \xmark & \xmark & \cmark & \pmark & \pmark & \pmark & \xmark \\

GuardAgent~\cite{xiang2025guardagent}
& \pmark & \pmark & \xmark & \xmark & \pmark & \pmark & \cmark & \pmark & \xmark \\

ShieldAgent~\cite{chen2025shieldagent}
& \pmark & \pmark & \xmark & \xmark & \cmark & \pmark & \cmark & \pmark & \xmark \\

AGrail~\cite{luo2025agrail}
& \pmark & \pmark & \xmark & \xmark & \pmark & \pmark & \cmark & \pmark & \xmark \\

\midrule
\textbf{\textsc{\proposed}}
& \cmark & \cmark & \cmark & \cmark & \cmark & \cmark & \cmark & \cmark & \cmark \\

\bottomrule
\end{tabular}
% \vspace{-0.3in}
\end{table*}

\paragraph{Agent Security Benchmarks:}
Agent security benchmarks have rapidly expanded from single-agent prompt-injection and tool-use settings to richer agentic environments. \textbf{Agent Security Bench (ASB)}~\cite{zhang2025agent} formalizes attacks and defenses for LLM-based agents across scenarios, tools, attack classes, and metrics, covering 10 application domains (e-commerce, autonomous driving, finance) with over 400 tools. \textbf{AgentDojo}~\cite{debenedetti2024agentdojo} evaluates tool-calling agents under prompt injection through untrusted data, providing realistic tasks and security test cases in an extensible environment spanning banking, travel, workspace, and Slack domains. \textbf{AgentHarm}~\cite{andriushchenko2025agentharm} measures whether agents comply with malicious requests and whether jailbreaks preserve multi-step task capability, using 110 harmful and 110 benign counterparts with 104 simulated tools. \textbf{SEC-bench}~\cite{lee2026secbench} evaluates LLM agents on realistic software-security tasks such as proof-of-concept generation and vulnerability patching, employing a multi-agent scaffold (SECVERIFIER) to automatically reproduce vulnerabilities and validate patches. These benchmarks primarily answer whether an attack succeeds or whether a defense blocks it; they do \textit{not} measure how much harm is produced per unit of local perturbation, nor do they systematically vary the size or subtlety of perturbations to quantify amplification. \proposed~ differs by introducing harm amplification, $\text{HA}(K)$, and by recording full execution traces, including specialist outputs, tool calls, shared memory, and oracle logs, before scoring. This enables reproducible measurement of how bounded local deviations cascade through systems with multiple agents and produce disproportionate downstream harm.

\paragraph{Multi-Agent Security:}
Multi-agent security work is conceptually closer to our setting. \textbf{TAMAS}~\cite{kavathekar2025tamas} benchmarks adversarial risks in multi-agent systems and includes prompt-level, environment-level, and agent-level attacks, covering scenarios with compromised and colluding agents. \textbf{Who is the Mole?}~\cite{xie2025s} studies intention-hiding malicious agents that behave plausibly while degrading system performance. Memory poisoning research~\cite{sunil2026memory} demonstrates that agents with persistent memory can continue harmful behavior even after the corrupting source disappears. These works introduce important attack vectors, including collusion, hidden intent, and memory persistence. However, they primarily measure attack \textit{success} rather than harm \textit{amplification}. TAMAS and Who is the Mole? do not quantify how final harm scales with the number or magnitude of compromised agents; they report whether the attack achieves its objective, not whether a small local deviation produces disproportionate downstream impact. Memory-poisoning work shows persistence but does not measure the ratio of downstream harm to initial perturbation size. \proposed~ extends this line of research by formalizing $\text{HA}(K)$ to capture how orchestration across multiple agents \textit{amplifies} local perturbations. It also separates explicit adversarial agents, such as single compromise and collusion, from correlated benign failure caused by shared biased context. This distinction is critical for understanding whether harm arises from intentional subversion or from systemic fragility in evidence aggregation.

\paragraph{Guardrail Systems:}
Safety guardrails for LLMs and agents aim to enforce security without modifying the foundation model. \textbf{LlamaGuard}~\cite{inan2023llama} pioneers LLM-based content moderation by fine-tuning models to classify inputs and outputs within a customizable safety framework. \textbf{Qwen3Guard}~\cite{zhao2025qwen3guard} introduces three-level harmfulness classification across 119 languages. \textbf{ShieldGemma}~\cite{zeng2024shieldgemma}, \textbf{PolyGuard}~\cite{kumar2025polyguard}, and \textbf{WildGuard}~\cite{han2024wildguard} provide additional content-moderation guardrails for static input and output. These models excel at filtering harmful prompts and responses but are \textit{not} designed to reason about dynamic multi-step agent trajectories or tool invocations. Recent agent-focused guardrails address this gap: \textbf{LlamaFirewall}~\cite{chennabasappa2025llamafirewall} combines PromptGuard2 with an alignment-check module to detect prompt injection and goal hijacking; \textbf{AgentAuditor}~\cite{luo2026agentauditor} retrieves reasoning experiences to guide LLM evaluation of complete execution trajectories. \textbf{ToolSafe}~\cite{mou2026toolsafe} introduces fine-grained safety checks for tool usage at each step through TS-Guard (a guardrail model trained with multi-task RL) and TS-Flow (a feedback-driven reasoning framework), reducing harmful tool invocations by 65\% while improving benign task completion by 10\% under prompt injection. \textbf{GuardAgent}~\cite{xiang2025guardagent} relies on user-specified guard rules, which can limit coverage. \textbf{ShieldAgent}~\cite{chen2025shieldagent} and \textbf{AGrail}~\cite{luo2025agrail} use richer multi-stage safety reasoning, which may add overhead, but their latency should be empirically reported rather than asserted.

\proposed~ complements these guardrail systems by providing a \textit{measurement framework} to evaluate their effectiveness specifically on harm amplification. Our experiments compare five defense settings, namely no defense, prompt sandwiching, LlamaFirewall, ToolSafe, and IntegrityGuard. We evaluate them not only on attack success rate, but also on harm amplification, benign utility, and stealth factor. This measurement capability is unique to \proposed~ and enables principled design choices that balance security, utility, and amplification mitigation.

\paragraph{Software Engineering Benchmarks:}
Software engineering (SE) benchmarks provide a useful contrast. \textbf{SWE-bench}~\cite{jimenez2024swebench} and its variants~\cite{yang2024sweagent} evaluate LLM agents on real-world bug-fixing issues from GitHub repositories. \textbf{Multi-SWE-bench}~\cite{zan2026multiswebench} and \textbf{SWE-PolyBench}~\cite{rashid2025swe} extend coverage to multiple programming languages. Additional benchmarks include \textbf{HumanEval}~\cite{chen2021evaluating}, \textbf{BigCodeBench}~\cite{zhuo2025bigcodebench}, \textbf{LiveCodeBench}~\cite{jain2025livecodebench}, and \textbf{EvalPlus}~\cite{liu2023is}. These benchmarks focus on code generation, debugging, and functional correctness; they are not designed to measure security vulnerabilities or harm propagation. While SE benchmarks demonstrate that state-of-the-art agents can achieve over 60\% success on real-world issues~\cite{anthropic2025claude37}, they do not address the question central to \proposed: when an agent makes a \textit{small} mistake in a multi-agent orchestration, how much harm does that mistake produce downstream? The functional-correctness framing in SE benchmarks treats all failures as equally important; \proposed~ instead measures whether the \textit{magnitude} of harm is proportional to the magnitude of the local deviation, revealing system-level fragility invisible to binary pass/fail metrics.

\noindent\textbf{Novelty of \proposed:}
\proposed~formulates harm amplification through role perturbation as a multi-agent attack vector. It measures how small, plausible changes to roles, shared context, memory, or tool observations become harmful final decisions through orchestration. Unlike prior benchmarks, \proposed~quantifies harm per local deviation and exposes systemic fragility.
